\documentclass[../main.tex]{subfiles}
\begin{document}

\section{The prediction shift phenomenon on ViT}
\label{sec:shift}

\begin{claim}{STRONG}{F09}
The published explanation of PSBD predicts that a clean input whose prediction moves under the perturbation lands on the attacker's target class, because the neuron bias the poisoning built favours that class once ordinary evidence is removed. The prediction is falsifiable on every cached cell. If the share of shifted clean predictions landing on the target does not exceed the benign reference wherever detection works, the neuron bias account is not what carries the method.
\end{claim}

\subsection{The measurement}

For every cached cell we take the recommended placement at the matched rule and at the adaptive rule, read the histogram of the classes that shifted clean predictions landed on, and compute the share on the target class, against the uniform expectation and against the same measurement on the matched dataset's benign reference probed with the same trigger, since a benign ViT already falls back onto a favourite class under perturbation and the panel's target is the same class \PanelTargetLabel{} everywhere. \Cref{tab:mech-shift-target} gives the share per attack and \cref{fig:shift-target} plots it against the achieved shift ratio, with the benign band. Over all cells the share is \ShiftToTargetShareAll{} against a benign share of \ShiftToTargetBenignShare{}, an excess of \ShiftToTargetExcessAll{}. The phenomenon exists on ViT.

% generated by scripts/paper/mech_shift_target.py from results/<folder>/psbd_metrics.json (65 clearing cells, 4 benign) at 2026-09-23T16:08:34.373666+00:00, commit afe5fd94d50b77f30f654fe075b768040987adef-dirty. Do not edit by hand.
\begin{table}[htbp]
\centering
\small
\caption{A1: target-class share of shifted clean predictions at the token\_mask placement at the attention input (before_attention_norm_token_mask), against the uniform expectation 1/num\_classes and the matched-dataset benign reference, read at the matched-0.6 and adaptive-0.8 rate rules. Excess is the attack's share minus the benign share on the same dataset, paired per model, 95\% bootstrap interval over models.}
\label{tab:mech-shift-target}
\begin{adjustbox}{max width=\linewidth}
\begin{tabular}{lrrrrrrr}
\toprule
attack & n & share@matched 0.6 & benign@matched 0.6 & excess@matched 0.6 [95\% CI] & share@adaptive 0.8 & benign@adaptive 0.8 & excess@adaptive 0.8 [95\% CI] \\
\midrule
BPP & 12 & 0.134 & 0.015 & +0.120 [+0.017, +0.250] & 0.168 & 0.001 & +0.166 [+0.039, +0.331] \\
WaNet & 5 & 0.339 & 0.023 & +0.315 [+0.076, +0.637] & 0.428 & 0.002 & +0.425 [+0.149, +0.715] \\
TaCT & 11 & 0.004 & 0.016 & $-$0.011 [$-$0.027, +0.003] & 0.002 & 0.002 & +0.000 [$-$0.002, +0.003] \\
SIG & 1 & 0.002 & 0.055 & $-$0.053 -- & 0.000 & 0.004 & $-$0.004 -- \\
BadNets & 12 & 0.276 & 0.015 & +0.261 [+0.080, +0.471] & 0.366 & 0.001 & +0.364 [+0.148, +0.593] \\
Blend & 12 & 0.015 & 0.015 & +0.000 [$-$0.014, +0.014] & 0.008 & 0.001 & +0.007 [+0.001, +0.015] \\
LF & 12 & 0.154 & 0.015 & +0.139 [+0.090, +0.191] & 0.217 & 0.001 & +0.215 [+0.132, +0.304] \\
all models & 65 & 0.134 & 0.016 & +0.118 [+0.066, +0.179] & 0.173 & 0.002 & +0.172 [+0.108, +0.242] \\
\bottomrule
\end{tabular}
\end{adjustbox}
\end{table}


\begin{figure}[t]
\centering
\includegraphics[width=0.7\linewidth]{figures/mech_shift_target.pdf}
\caption{Target-class share of shifted clean predictions against the achieved clean shift ratio, a line per attack averaged over its cached cells, the benign references as a band and the uniform expectation dashed. From \code{paper/figures/mech\_shift\_target.json}.}
\label{fig:shift-target}
\end{figure}

\subsection{Where it exists and where detection works}

\Cref{tab:mech-shift-target-by-dataset} is the honest version of the table, because the per-attack means hide that the effect is a property of the dataset as much as of the attack. The GTSRB BadNets and WaNet cells put most shifted clean predictions on the target, the CIFAR-10 BadNets cells a third, the BPP and low-frequency cells a modest share everywhere, and the CIFAR-100 BadNets cells, Blend on every dataset and TaCT on every dataset put none there beyond the benign reference. On CIFAR-100 BadNets at \PanelRateLow{} the share at the adaptive rule is \ShiftToTargetHeadlineCell{}, against the uniform expectation of a single class in \ClassesCifarOneZeroZero{}, and that cell is where the recommended placement gains most over the published one. The phenomenon is a local-trigger, few-class effect. A patch trigger on a dataset with tens of classes drags the fallback class onto the target, a global trigger does not, and the number of classes dilutes it, which is the pattern Li et al. themselves saw on Tiny ImageNet and attributed to generalisation.

% generated by scripts/paper/mech_shift_target.py from results/<folder>/psbd_metrics.json (65 clearing cells, 4 benign) at 2026-09-23T16:54:56.873920+00:00, commit d4a8014354b8fcc8d014beec50fef60ed46fa47d-dirty. Do not edit by hand.
\begin{table}[htbp]
\centering
\small
\caption{Share of shifted clean predictions that land on the target class under PSBD-TM, per dataset and attack, against the uniform expectation and the benign reference of the dataset.}
\label{tab:mech-shift-target-by-dataset}
\begin{adjustbox}{max width=\linewidth}
\begin{tabular}{lrrrrrr}
\toprule
Dataset & Attack & n & Share & Uniform & Benign & Excess \\
\midrule
CIFAR-10 & BPP & 3 & 0.166 & 0.100 & 0.004 & +0.162 \\
CIFAR-10 & WaNet & 2 & 0.172 & 0.100 & 0.004 & +0.168 \\
CIFAR-10 & TaCT & 3 & 0.000 & 0.100 & 0.004 & $-$0.004 \\
CIFAR-10 & SIG & 1 & 0.000 & 0.100 & 0.004 & $-$0.004 \\
CIFAR-10 & BadNets & 3 & 0.518 & 0.100 & 0.004 & +0.514 \\
CIFAR-10 & Blend & 3 & 0.005 & 0.100 & 0.004 & +0.001 \\
CIFAR-10 & LF & 3 & 0.172 & 0.100 & 0.004 & +0.168 \\
CIFAR-100 & BPP & 3 & 0.157 & 0.010 & 0.000 & +0.157 \\
CIFAR-100 & TaCT & 3 & 0.004 & 0.010 & 0.000 & +0.004 \\
CIFAR-100 & BadNets & 3 & 0.002 & 0.010 & 0.000 & +0.002 \\
CIFAR-100 & Blend & 3 & 0.009 & 0.010 & 0.000 & +0.009 \\
CIFAR-100 & LF & 3 & 0.256 & 0.010 & 0.000 & +0.256 \\
GTSRB & BPP & 3 & 0.001 & 0.023 & 0.000 & +0.001 \\
GTSRB & WaNet & 1 & 0.996 & 0.023 & 0.000 & +0.996 \\
GTSRB & TaCT & 2 & 0.000 & 0.023 & 0.000 & +0.000 \\
GTSRB & BadNets & 3 & 0.793 & 0.023 & 0.000 & +0.793 \\
GTSRB & Blend & 3 & 0.015 & 0.023 & 0.000 & +0.015 \\
GTSRB & LF & 3 & 0.054 & 0.023 & 0.000 & +0.054 \\
Tiny ImageNet & BPP & 3 & 0.347 & 0.005 & 0.002 & +0.346 \\
Tiny ImageNet & WaNet & 2 & 0.400 & 0.005 & 0.002 & +0.398 \\
Tiny ImageNet & TaCT & 3 & 0.002 & 0.005 & 0.002 & +0.000 \\
Tiny ImageNet & BadNets & 3 & 0.150 & 0.005 & 0.002 & +0.149 \\
Tiny ImageNet & Blend & 3 & 0.005 & 0.005 & 0.002 & +0.003 \\
Tiny ImageNet & LF & 3 & 0.385 & 0.005 & 0.002 & +0.383 \\
\bottomrule
\end{tabular}
\end{adjustbox}
\end{table}


Detection does not follow the phenomenon. Blend is among the best-detected attacks on every dataset and shows no shift to the target anywhere, CIFAR-100 BadNets is detected at a mean AUROC above the panel mean and shows none, and the benign references show a share near the uniform line while detection sits at chance. The concentration half of the paper's observation holds, since shifted predictions do collapse onto a single class on ViT, but that class is a property of the model and of the trigger geometry rather than of the attacker's target. The neuron bias account is a correct description of some cells and it is not what carries the method.

\subsection{Removal against disturbance, and the confidence null}

If the mechanism were removal of the neurons that carry clean evidence, an operator that removes nothing should fail. Gaussian noise removes nothing and, at the normalised MLP input where no LayerNorm absorbs it, beats the token mask by \GaussianMinusTokenMaskMlp{} in \cref{tab:operators-deltas}, so removal is not required\finding{STRONG}{F06}. If instead the perturbation earned nothing at all and PSU measured the unperturbed confidence, the max-softmax confidence would match it. \Cref{tab:confidence-null} scores both on the same cached tensors over \ConfidenceNullCheckpoints{} ViT checkpoints at the residual placement the record was written for, and PSU reads \ConfidenceNullPsuAuroc{} against \ConfidenceNullAuroc{} for the confidence null, winning on \ConfidenceNullPsuWins{} of them\finding{SUPPORTING}{F10}. The perturbation does real work and it does not need to remove anything, which leaves an account in which PSU measures how far a prediction sits from a decision boundary along whatever direction the operator's Jacobian selects, a triggered input sitting far from it because the trigger's direction was fitted faster and more completely than any clean feature. That account is stated as a measurement of what the method needs and does not need, and \cref{sec:planned} lists the direct test of the margin reading that has not been run.

% generated by scripts/paper/mech_confidence_null.py from results/_experiments/psu_vs_confidence/psu_vs_confidence.json at 2026-09-23T16:50:43.186997+00:00, commit d4a8014354b8fcc8d014beec50fef60ed46fa47d-dirty. Do not edit by hand.
\begin{table}[htbp]
\centering
\small
\caption{PSU against the confidence null at the `pre_residual` placement: mean AUROC of PSU and of the unperturbed max-softmax confidence on the same cached tensors, over every ViT backdoored checkpoint in the record outside the SAM, evasion and seed sets, and the count of checkpoints where PSU wins.}
\label{tab:confidence-null}
\begin{adjustbox}{max width=\linewidth}
\begin{tabular}{lrrrrr}
\toprule
dataset & n & PSU AUROC & confidence AUROC & PSU minus confidence & PSU wins \\
\midrule
CIFAR-10 & 8 & 0.902 & 0.644 & +0.258 & 8 \\
CIFAR-100 & 7 & 0.833 & 0.757 & +0.076 & 4 \\
GTSRB & 6 & 0.990 & 0.372 & +0.618 & 6 \\
Tiny ImageNet & 7 & 0.836 & 0.819 & +0.017 & 2 \\
all & 28 & 0.887 & 0.658 & +0.229 & 20 \\
\bottomrule
\end{tabular}
\end{adjustbox}
\end{table}


\end{document}
